\subsection{Finite Sums}

\begin{definition} \label{def:finite-sums-indexed-by-integers} \hlt{Finite Sums Indexed by Integers}\\
Let $\mathbb{F}$ be a field and let $(a_i)_{i\in\Z}$ be a family in $\mathbb{F}$. Fix $m,n\in\Z$. If $m\leq n$, define $\sum\limits_{i=m}^{n}a_i\defeq a_m+a_{m+1}+\cdots+a_n$.\\
If $m>n$, set $\sum\limits_{i=m}^{n}a_i\defeq0$; this is the empty sum.
\end{definition}

\begin{proposition} \label{prop:finite-sum-basic-laws} \hlt{Basic Laws for Finite Sums}\\
Let $\mathbb{F}$ be a field. Let $(a_i)_{i\in\Z}$ and $(b_i)_{i\in\Z}$ be families in $\mathbb{F}$. Fix $m,n\in\Z$ with $m\leq n$.
\begin{enumerate}[label=(\roman*)]
\setlength{\itemsep}{0pt}
\item \hlt{Splitting}: Fix $r\in\Z$ with $m\leq r<n$. Then $\sum\limits_{i=m}^{n}a_i=\sum\limits_{i=m}^{r}a_i+\sum\limits_{i=r+1}^{n}a_i$.
\item \hlt{Additivity}: $\sum\limits_{i=m}^{n}(a_i+b_i)=\sum\limits_{i=m}^{n}a_i+\sum\limits_{i=m}^{n}b_i$.
\item \hlt{Homogeneity}: For $c\in\mathbb{F}$, $\sum\limits_{i=m}^{n}ca_i=c\sum\limits_{i=m}^{n}a_i$.
\item \hlt{Triangle inequality}: If $\mathbb{F}=\C$, then $\abs{\sum\limits_{i=m}^{n}a_i}\leq\sum\limits_{i=m}^{n}\abs{a_i}$.
\item \hlt{Monotonicity}: Suppose $\mathbb{F}=\R$ and $(\forall i\in\Z)\,[m\leq i\leq n\Rightarrow a_i\leq b_i]$. Then $\sum\limits_{i=m}^{n}a_i\leq\sum\limits_{i=m}^{n}b_i$.
\end{enumerate}
\begin{proof} For (ii)--(v), fix $m$ and set $\ell\defeq n-m+1$; by Theorem \ref{thm:mathematical-induction} with induction parameter $\ell$, it suffices in each item to verify the case $n=m$ and the step from $n$ to $n+1$.
\begin{enumerate}[label=(\roman*)]
\setlength{\itemsep}{0pt}
\item By Definition \ref{def:finite-sums-indexed-by-integers} and associativity of addition in $\mathbb{F}$, $\sum\limits_{i=m}^{n}a_i=\sum\limits_{i=m}^{r}a_i+\sum\limits_{i=r+1}^{n}a_i$.
\item If $n=m$, both sides equal $a_m+b_m$. Assume additivity through $n$. By (i), $\sum\limits_{i=m}^{n+1}(a_i+b_i)=\sum\limits_{i=m}^{n}(a_i+b_i)+a_{n+1}+b_{n+1}$.\\
Substitute the induction hypothesis and regroup to obtain $\sum\limits_{i=m}^{n+1}a_i+\sum\limits_{i=m}^{n+1}b_i$.
\item If $n=m$, both sides equal $ca_m$. Assume homogeneity through $n$. By (i), $\sum\limits_{i=m}^{n+1}ca_i=\sum\limits_{i=m}^{n}ca_i+ca_{n+1}$.\\
The induction hypothesis and distributivity turn the right side into $c(\sum\limits_{i=m}^{n}a_i+a_{n+1})$. By (i), this equals $c\sum\limits_{i=m}^{n+1}a_i$.
\item If $n=m$, equality holds. Assume the inequality through $n$. By (i) and the triangle inequality, $\abs{\sum\limits_{i=m}^{n+1}a_i}\leq\abs{\sum\limits_{i=m}^{n}a_i}+\abs{a_{n+1}}$.\\
The induction hypothesis and (i) bound the right side by $\sum\limits_{i=m}^{n+1}\abs{a_i}$.
\item If $n=m$, the conclusion is the hypothesis $a_m\leq b_m$. Assume monotonicity through $n$. By (i), the sums through $n+1$ split at $n$.\\
The induction hypothesis, order compatibility, and $a_{n+1}\leq b_{n+1}$ give $\sum\limits_{i=m}^{n+1}a_i\leq\sum\limits_{i=m}^{n+1}b_i$.
\end{enumerate}
\hlt{Structural Commentary}: Splitting reduces each law to one new summand and the scalar axioms.\\
\hlt{Key Insight}: Finite-sum laws propagate one term at a time.
\end{proof}
\end{proposition}
